% GENERATED FILE. Edit canonical lessons, then run npm run generate:books.
% canonical-source-hash: fnv1a64:244422db3a15375f
% canonical-lessons: BN-C519-gelari, BN-C519-mech, BN-C519-odhinayok, BN-C519-rephari, BN-C519-hara

\chapter{The Match}
\label{ch:bn-the-match}
\hlchaptermodality{519}

\begin{chapteropening}
\textbf{By the end of this chapter:} \emph{I can talk about a match, its captain, its referee and the stands.}

Complete the last lesson of chapter 519: I can talk about a match, its captain, its referee and the stands.
\end{chapteropening}

\section[\texorpdfstring{\bn{গ্যালারি} (gælāri)}{গ্যালারি (gælāri)}]{\bn{গ্যালারি} (gælāri) — the stands (at a stadium); a gallery}

\label{lesson:BN-C519-gelari}



\begin{quote}
Before the new one: say the Bengali for a crossword, then the Bengali for carrom.
\end{quote}

\subsection*{You'll want to know: \bn{গ্যালারি}}
\textbf{\bn{গ্যালারি}} — \emph{gælāri} — \textquotedblleft{}the stands (at a stadium); a gallery\textquotedblright{}.

\begin{quote}
\bn{স্টেডিয়ামের} \bn{গ্যালারি} \bn{দর্শকে} \bn{ভর্তি} \bn{ছিল।}
\end{quote}

\begin{quote}
\emph{Sṭeḍiyāmer gælāri dôrshôke bhôrti chhilô.}
\end{quote}

\begin{quote}
The stands at the stadium were full of spectators.
\end{quote}

\begin{cousinweb}
\textbf{\bn{গ্যালারি}} (\emph{gælāri}) is borrowed from English \emph{gallery}, used for the stands around a ground.
\end{cousinweb}

\begin{grammarlens}[title={What you've built}]
One more word for describing what you lived and felt.
\end{grammarlens}

\subsection*{Your turn}
\begin{itemize}
\raggedright
  \item \emph{Say it:} \emph{gælāri}
  \item \emph{Say it:} \emph{gælāri}, once more
  \item \emph{Say it:} \emph{Sṭeḍiyāmer gælāri dôrshôke bhôrti chhilô.} — The stands at the stadium were full of spectators.
  \item \emph{Recall:} say the Bengali for a crossword, then the Bengali for carrom, then say \emph{gælāri} again
\end{itemize}

\begin{tcolorbox}[breakable,colback=teal!4,colframe=teal!35!black,title={Before you move on}]
What does \bn{গ্যালারি} mean? (\textquotedblleft{}The stands (at a stadium); a gallery\textquotedblright{}.) Say it once more.
\end{tcolorbox}
\section[\texorpdfstring{\bn{ম্যাচ} (mæch)}{ম্যাচ (mæch)}]{\bn{ম্যাচ} (mæch) — a match (a game)}

\label{lesson:BN-C519-mech}



\begin{quote}
Before the new one: say the Bengali for carrom, then the Bengali for the stands.
\end{quote}

\subsection*{You'll want to know: \bn{ম্যাচ}}
\textbf{\bn{ম্যাচ}} — \emph{mæch} — \textquotedblleft{}a match (a game)\textquotedblright{}.

\begin{quote}
\bn{রবিবার} \bn{আমাদের} \bn{ফুটবল} \bn{ম্যাচ।}
\end{quote}

\begin{quote}
\emph{Rôbibār āmāder phuṭbôl mæch.}
\end{quote}

\begin{quote}
We have a football match on Sunday.
\end{quote}

\begin{cousinweb}
\textbf{\bn{ম্যাচ}} (\emph{mæch}) is borrowed from English \emph{match}.
\end{cousinweb}

\begin{grammarlens}[title={What you've built}]
One more word for describing what you lived and felt.
\end{grammarlens}

\subsection*{Your turn}
\begin{itemize}
\raggedright
  \item \emph{Say it:} \emph{mæch}
  \item \emph{Say it:} \emph{mæch}, once more
  \item \emph{Say it:} \emph{Rôbibār āmāder phuṭbôl mæch.} — We have a football match on Sunday.
  \item \emph{Recall:} say the Bengali for carrom, then the Bengali for the stands, then say \emph{mæch} again
\end{itemize}

\begin{tcolorbox}[breakable,colback=teal!4,colframe=teal!35!black,title={Before you move on}]
What does \bn{ম্যাচ} mean? (\textquotedblleft{}A match (a game)\textquotedblright{}.) Say it once more.
\end{tcolorbox}
\section[\texorpdfstring{\bn{অধিনায়ক} (ôdhināyôk)}{অধিনায়ক (ôdhināyôk)}]{\bn{অধিনায়ক} (ôdhināyôk) — a team captain}

\label{lesson:BN-C519-odhinayok}



\begin{quote}
Before the new one: say the Bengali for the stands, then the Bengali for a match.
\end{quote}

\subsection*{You'll want to know: \bn{অধিনায়ক}}
\textbf{\bn{অধিনায়ক}} — \emph{ôdhināyôk} — \textquotedblleft{}a team captain\textquotedblright{}.

\begin{quote}
\bn{আমার} \bn{বন্ধু} \bn{দলের} \bn{অধিনায়ক।}
\end{quote}

\begin{quote}
\emph{Āmār bôndhu dôler ôdhināyôk.}
\end{quote}

\begin{quote}
My friend is the captain of the team.
\end{quote}

\begin{cousinweb}
\textbf{\bn{অধিনায়ক}} (\emph{ôdhināyôk}) joins Sanskrit \emph{adhi}, \textquotedblleft{}over, above\textquotedblright{}, to \emph{nāyaka}, \textquotedblleft{}a guide, a leader\textquotedblright{} (Monier-Williams): the leader set over the rest.
\end{cousinweb}

\begin{grammarlens}[title={What you've built}]
One more word for describing what you lived and felt.
\end{grammarlens}

\subsection*{Your turn}
\begin{itemize}
\raggedright
  \item \emph{Say it:} \emph{ôdhināyôk}
  \item \emph{Say it:} \emph{ôdhināyôk}, once more
  \item \emph{Say it:} \emph{Āmār bôndhu dôler ôdhināyôk.} — My friend is the captain of the team.
  \item \emph{Recall:} say the Bengali for the stands, then the Bengali for a match, then say \emph{ôdhināyôk} again
\end{itemize}

\begin{tcolorbox}[breakable,colback=teal!4,colframe=teal!35!black,title={Before you move on}]
What does \bn{অধিনায়ক} mean? (\textquotedblleft{}A team captain\textquotedblright{}.) Say it once more.
\end{tcolorbox}
\section[\texorpdfstring{\bn{রেফারি} (rephāri)}{রেফারি (rephāri)}]{\bn{রেফারি} (rephāri) — a referee}

\label{lesson:BN-C519-rephari}



\begin{quote}
Before the new one: say the Bengali for a match, then the Bengali for a team captain.
\end{quote}

\subsection*{You'll want to know: \bn{রেফারি}}
\textbf{\bn{রেফারি}} — \emph{rephāri} — \textquotedblleft{}a referee\textquotedblright{}.

\begin{quote}
\bn{রেফারি} \bn{বাঁশি} \bn{বাজালেন।}
\end{quote}

\begin{quote}
\emph{Rephāri bā̃shi bājālen.}
\end{quote}

\begin{quote}
The referee blew the whistle.
\end{quote}

\begin{cousinweb}
\textbf{\bn{রেফারি}} (\emph{rephāri}) is borrowed from English \emph{referee}.
\end{cousinweb}

\begin{grammarlens}[title={What you've built}]
One more word for describing what you lived and felt.
\end{grammarlens}

\subsection*{Your turn}
\begin{itemize}
\raggedright
  \item \emph{Say it:} \emph{rephāri}
  \item \emph{Say it:} \emph{rephāri}, once more
  \item \emph{Say it:} \emph{Rephāri bā̃shi bājālen.} — The referee blew the whistle.
  \item \emph{Recall:} say the Bengali for a match, then the Bengali for a team captain, then say \emph{rephāri} again
\end{itemize}

\begin{tcolorbox}[breakable,colback=teal!4,colframe=teal!35!black,title={Before you move on}]
What does \bn{রেফারি} mean? (\textquotedblleft{}A referee\textquotedblright{}.) Say it once more.
\end{tcolorbox}
\section[\texorpdfstring{\bn{হারা} (hārā)}{হারা (hārā)}]{\bn{হারা} (hārā) — to lose (a game); to be beaten}

\label{lesson:BN-C519-hara}



\begin{quote}
Before the new one: say the Bengali for a team captain, then the Bengali for a referee.
\end{quote}

\subsection*{You'll want to know: \bn{হারা}}
\textbf{\bn{হারা}} — \emph{hārā} — \textquotedblleft{}to lose (a game); to be beaten\textquotedblright{}.

\begin{quote}
\bn{কাল} \bn{আমাদের} \bn{দল} \bn{হারল।}
\end{quote}

\begin{quote}
\emph{Kāl āmāder dôl hārlô.}
\end{quote}

\begin{quote}
Our team lost yesterday.
\end{quote}

Past (he/she …): \bn{তিনি} \bn{দাবায়} \bn{হারলেন}, \emph{tini dābāy hārlen}.

\begin{cousinweb}
\textbf{\bn{হারা}} (\emph{hārā}) is the plain partner of \bn{হারানো}, \textquotedblleft{}to lose (a thing)\textquotedblright{}; CDIAL lists Bengali \emph{hārā}, \textquotedblleft{}to lose, to be defeated\textquotedblright{} (CDIAL 14061).
\end{cousinweb}

\begin{grammarlens}[title={What you've built}]
One more word for describing what you lived and felt.
\end{grammarlens}

\subsection*{Your turn}
\begin{itemize}
\raggedright
  \item \emph{Say it:} \emph{hārā}
  \item \emph{Say it:} \emph{hārā}, once more
  \item \emph{Say it:} \emph{Kāl āmāder dôl hārlô.} — Our team lost yesterday.
  \item \emph{Recall:} say the Bengali for a team captain, then the Bengali for a referee, then say \emph{hārā} again
\end{itemize}

\begin{tcolorbox}[breakable,colback=teal!4,colframe=teal!35!black,title={Before you move on}]
What does \bn{হারা} mean? (\textquotedblleft{}To lose (a game); to be beaten\textquotedblright{}.) Say it once more.
\end{tcolorbox}
